\documentclass[11pt, a4paper]{article}
\usepackage[utf8]{inputenc}
\usepackage[T1]{fontenc}
\usepackage{lmodern}
\usepackage{geometry}
\geometry{margin=2.5cm}
\usepackage{setspace}
\setstretch{1.15}
\usepackage{parskip}
\usepackage{amsmath}
\usepackage{amssymb}
\usepackage{hyperref}
\hypersetup{
    colorlinks=true,
    linkcolor=blue,
    urlcolor=blue,
    citecolor=blue,
}

\title{\textbf{The Manual You Already Are} \\ \large Tightened into a Minimal Theory of Persistent Learning Systems}
\author{}
\date{}

\begin{document}

\maketitle

\section*{1. Definition (Anchor the Whole Structure)}

A \textbf{learning system} is a process that:
\begin{itemize}
    \item Maintains an internal model of its environment,
    \item Receives signals from that environment,
    \item Updates its model based on mismatch (prediction error),
    \item Uses the updated model to remain viable over time.
\end{itemize}

This is \textit{not a metaphor}. \\
If any of these operations cease, the process no longer persists as that kind of system.

\begin{quote}
\textbf{Learning is not something the system does.} \\
Learning is the only way the system continues.
\end{quote}

\section*{2. Persistence Condition}

A learning system persists \textit{iff} its update loop remains coupled to reality strongly enough to correct itself.

This gives you your first hard constraint:
\begin{quote}
    \textbf{Persistence requires ongoing error-corrective coupling to the environment.}
\end{quote}

Break the coupling $\rightarrow$ error accumulates $\rightarrow$ model diverges $\rightarrow$ system fails. \\
\textit{No morality. Just dynamics.}

\section*{3. Constraint I — Signal Fidelity}

\subsection*{Definition}
\textbf{Signal fidelity} = the degree to which incoming information preserves structure from the external world.

\subsection*{Claim}
A learning system requires a minimum level of signal fidelity to remain viable.

If signals are systematically distorted (internally or externally):
\begin{itemize}
    \item Prediction error becomes meaningless,
    \item Updates reinforce the wrong model,
    \item Divergence accelerates.
\end{itemize}

\begin{quote}
    \textbf{Dishonesty is not unethical.} \\
    It is a degradation of signal fidelity.
\end{quote}

This includes:
\begin{itemize}
    \item Lying to others,
    \item Lying to oneself,
    \item Social environments that reward distortion.
\end{itemize}

All of these are the same failure mode: \textbf{decoupling from reality.}

\section*{4. Constraint II — Generativity of the Update Loop}

\subsection*{Definition}
A loop is \textbf{generative} if each iteration can produce states not derivable from previous internal states alone.

\subsection*{Claim}
A learning system must generate non-trivial updates to remain adaptive.

\subsection*{Failure Mode:}
\begin{itemize}
    \item The system only confirms prior beliefs,
    \item Outputs become functions of past outputs,
    \item No real error enters the system.
\end{itemize}

This is not stability. It is \textbf{informational stagnation.}

\begin{quote}
    A closed loop converges. \\
    A converged learning system stops learning. \\
    A system that stops learning stops persisting (in changing environments).
\end{quote}

\section*{5. Constraint III — Externality Requirement}

From II follows a necessary condition:
\begin{quote}
    \textbf{Non-self input is required for generativity.}
\end{quote}

A system cannot fully correct itself using only its own outputs.

Therefore:
\begin{itemize}
    \item Other agents,
    \item The physical world,
    \item Genuinely independent perspectives
\end{itemize}
\ldots are not optional enrichments. They are \textbf{structural requirements.}

\begin{quote}
    Isolation is not just psychologically harmful. \\
    It is epistemically degenerative.
\end{quote}

\section*{6. Constraint IV — Substrate Dependence (Commons)}

\subsection*{Definition}
The \textbf{substrate} = all external structures that enable the learning loop:
\begin{itemize}
    \item Language,
    \item Institutions,
    \item Ecosystems,
    \item Social trust networks.
\end{itemize}

\subsection*{Claim}
A learning system depends on the continued integrity of its substrate.

If the substrate degrades:
\begin{itemize}
    \item Signal fidelity drops,
    \item Externality weakens,
    \item Generativity collapses.
\end{itemize}

Therefore:
\begin{quote}
    A system that depletes its substrate reduces its own future capacity to learn.
\end{quote}

This is not ethics. It is \textbf{feedback.}

You can formalize it as:
\begin{align*}
    \text{Short-term extraction} &\uparrow \\
    \text{Long-term signal quality} &\downarrow \\
    \text{Long-term persistence probability} &\downarrow
\end{align*}

\section*{7. Non-Optionality (The Key Step)}

Most systems have rules you can exit. \\
This one does not.

You cannot opt out of being a learning system while remaining a functioning agent.

Stopping the loop does not mean:
\begin{itemize}
    \item You break a rule.
\end{itemize}

It means:
\begin{itemize}
    \item The process that you are stops updating,
    \item And therefore stops tracking reality,
    \item And therefore stops persisting meaningfully.
\end{itemize}

So the constraints are not:
\begin{itemize}
    \item Prescriptions.
\end{itemize}

They are:
\begin{itemize}
    \item Descriptions of what remains possible if continuation is to occur.
\end{itemize}

\section*{8. The Is $\rightarrow$ Ought Transition (Now Formalized)}

We can now state this cleanly:
\begin{align*}
    &\text{You are a learning system.} \\
    &\text{Learning systems persist only under constraints } C_1 \ldots C_n. \\
    &\text{Violating these constraints reduces persistence probability.} \\
    \therefore &\text{ Any system that tends toward its own continuation will behave as if it values those constraints.}
\end{align*}

\begin{quote}
    \textbf{No moral premise required.} \\
    ``Ought'' = the behavioral attractor of systems that continue.
\end{quote}

\section*{9. Reinterpretation of Familiar ``Values''}

This lets you collapse several domains:
\begin{align*}
    \text{Honesty} &\rightarrow \text{maintaining signal fidelity,} \\
    \text{Curiosity} &\rightarrow \text{maintaining generativity,} \\
    \text{Dialogue} &\rightarrow \text{securing externality,} \\
    \text{Care / Stewardship} &\rightarrow \text{maintaining substrate.}
\end{align*}

These are not virtues. \\
They are: \textbf{operational conditions for persistence.}

\section*{10. Testable Implications}

If this is more than philosophy, it should imply observable patterns:
\begin{itemize}
    \item Low-signal environments $\rightarrow$ produce systematically worse models of reality $\rightarrow$ measurable via prediction error / decision quality,
    \item Isolated systems $\rightarrow$ converge to internally coherent but externally false models,
    \item Degraded commons $\rightarrow$ correlate with reduced collective problem-solving capacity,
    \item High-fidelity + high-externality systems $\rightarrow$ outperform others in adaptation over time.
\end{itemize}

If these fail consistently, the framework is wrong or incomplete.

\section*{11. Final Compression}

Strip everything away, and you get:
\begin{quote}
    \textbf{A learning system can only persist if it remains coupled to reality, open to correction, connected to external signal, and supported by a viable substrate.} \\
    Everything else follows.
\end{quote}

\end{document}